\subsection{Recubrimientos étales.}
\label{subsection:IV.18.2}

\oldpage[IV-4]{111}

\begin{env}[18.2.1]
\label{IV.18.2.1}
Dados un anillo $A$ y una $A$-álgebra conmutativa $B$ que es \emph{finita} y es un $A$-módulo \emph{libre}, recordemos (Bourbaki, \emph{Alg.}, cap. VIII, \textsection~12, n.\textsuperscript{o}~2) que se define sobre $B$ una \emph{forma $A$-lineal} $\operatorname{Tr}_{B/A}$, la «forma traza»; de ella se deduce la definición de una \emph{forma $A$-bilineal simétrica} (también llamada «forma traza»)
\begin{equation}
  (x,y)\longmapsto\operatorname{Tr}_{B/A}(xy)
  \tag{18.2.1.1}\label{IV.18.2.1.1}
\end{equation}
cuyo dato equivale al de la aplicación $A$-lineal asociada $\operatorname{astr}_{B/A}:B\to\check B$ del $A$-módulo $B$ a su \emph{dual} $\check B$, igual a su traspuesta. Cuando $A$ es un \emph{cuerpo}, es equivalente decir que esta forma bilineal es \emph{no degenerada} o que $B$ es una $A$-álgebra \emph{separable} (Bourbaki, \emph{Alg.}, cap. IX, \textsection~2, prop. 5).

Sea $f:A\to A'$ un homomorfismo de anillos; pongamos $B'=B\otimes_A A'$, y sea $g:B\to B'$ el homomorfismo canónico; la imagen por $g$ de una base del $A$-módulo $B$ es entonces una base del $A'$-módulo $B'$, y se sigue de las definiciones que se tiene, para todo $x\in B$,
\begin{equation}
  \operatorname{Tr}_{B'/A'}(g(x))=f(\operatorname{Tr}_{B/A}(x)).
  \tag{18.2.1.2}\label{IV.18.2.1.2}
\end{equation}
\end{env}

\begin{env}[18.2.2]
\label{IV.18.2.2}
Consideremos ahora un \emph{espacio anillado} $(X,\sh{O}_X)$ y sea $\sh{B}$ una $\sh{O}_X$-álgebra que, como $\sh{O}_X$-módulo, es \emph{localmente libre de rango finito}; entonces, para todo abierto $U\subset X$ tal que $\sh{B}|U$ sea (como $\sh{O}_U$-módulo) isomorfa a $\sh{O}_U^n$ (para un $n$ que depende de $U$), $\Gamma(U,\sh{B})$ es una $\Gamma(U,\sh{O}_X)$-álgebra que, como $\Gamma(U,\sh{O}_X)$-módulo, es libre de rango finito y define, por tanto, una forma $\Gamma(U,\sh{O}_X)$-lineal $\operatorname{Tr}_{\Gamma(U,\sh{B})/\Gamma(U,\sh{O}_X)}$, que denotaremos también por $\operatorname{Tr}_{\sh{B}|\sh{O}_X,U}$; de ella se deduce una aplicación lineal asociada
\[
  \operatorname{astr}_{\sh{B}|\sh{O}_X,U}:\Gamma(U,\sh{B})\longrightarrow\Gamma(U,\sh{B})^\vee=\Gamma(U,\check{\sh{B}}).
\]
Además, se sigue de (18.2.1.2) que estas aplicaciones lineales son compatibles con las operaciones de restricción de $U$ a un abierto menor y definen, por una parte, un \emph{homomorfismo de $\sh{O}_X$-módulos}, llamado también \emph{homomorfismo traza}:
\begin{equation}
  \operatorname{Tr}_{\sh{B}|\sh{O}_X}:\sh{B}\longrightarrow\sh{O}_X
  \tag{18.2.2.1}\label{IV.18.2.2.1}
\end{equation}
y, por otra parte, un \emph{homomorfismo de $\sh{O}_X$-módulos}
\begin{equation}
  \operatorname{astr}_{\sh{B}|\sh{O}_X}:\sh{B}\longrightarrow\check{\sh{B}}
  \tag{18.2.2.2}\label{IV.18.2.2.2}
\end{equation}
llamado \emph{asociado a la traza}, e igual a su traspuesto. Se sigue también de (18.2.1.2) que, para todo $x\in X$, se tiene
\begin{equation}
  (\operatorname{Tr}_{\sh{B}|\sh{O}_X})_x=\operatorname{Tr}_{\sh{B}_x|\sh{O}_{X,x}}
  \tag{18.2.2.3}\label{IV.18.2.2.3}
\end{equation}
\begin{equation}
  (\operatorname{astr}_{\sh{B}|\sh{O}_X})_x=\operatorname{astr}_{\sh{B}_x|\sh{O}_{X,x}}.
  \tag{18.2.2.4}\label{IV.18.2.2.4}
\end{equation}
Finalmente, bajo las condiciones de (18.2.1), si se pone $X=\Spec(A)$ y si $\sh{B}=\widetilde B$ es la $\sh{O}_X$-álgebra correspondiente a $B$, la forma $\operatorname{Tr}_{\sh{B}|\sh{O}_X}$ (resp. el homomorfismo $\operatorname{astr}_{\sh{B}|\sh{O}_X}$) corresponde a la forma $\operatorname{Tr}_{B/A}$ (resp. al homomorfismo de $A$-módulos $\operatorname{astr}_{B/A}$), como se sigue igualmente de (18.2.1.2).
\end{env}


\oldpage[IV-4]{112}

\begin{proposition}[18.2.3]
\label{IV.18.2.3}
Sea $f:X\to Y$ un morfismo finito de preesquemas y sea $\sh{B}=f_*(\sh{O}_X)$. Las condiciones siguientes son equivalentes:
\begin{enumerate}
  \item[a)] $f$ es étale.
  \item[a$'$)] $f$ es un morfismo plano de presentación finita, y para todo $x\in X$, si se pone $y=f(x)$, $\sh{O}_{X,x}/\mathfrak m_y\sh{O}_{X,x}$ es un cuerpo, extensión finita separable de $\kappa(y)$.
  \item[b)] $\sh{B}$ es un $\sh{O}_Y$-módulo localmente libre y, para todo $y\in Y$, $\sh{B}_y\otimes_{\sh{O}_{Y,y}}\kappa(y)$ es una $\kappa(y)$-álgebra finita separable (por tanto, el producto directo de un número finito de cuerpos, extensiones finitas separables de $\kappa(y)$).
  \item[c)] $\sh{B}$ es un $\sh{O}_Y$-módulo localmente libre, y el homomorfismo $\operatorname{astr}_{\sh{B}|\sh{O}_Y}:\sh{B}\to\check{\sh{B}}$ (18.2.2) es biyectivo.
\end{enumerate}
\end{proposition}

Habida cuenta de que $f$ es cuasicompacto, la equivalencia de $a)$ y $a')$ ya ha sido demostrada (17.6.2). Para demostrar el resto de la proposición, se puede reducir al caso en que $Y=\Spec(A)$ y $X=\Spec(B)$ son afines, siendo $B$ una $A$-álgebra finita y $\sh{B}=\widetilde B$. Decir que $f$ es un morfismo de presentación finita equivale entonces a decir que $B$ es un $A$-módulo de presentación finita (1.4.7). Si además $f$ es plano, y por tanto $B$ un $A$-módulo plano, se sabe (Bourbaki, \emph{Alg. comm.}, cap. II, \textsection~5, n.\textsuperscript{o}~2, cor. 2 del th. 1) que $B$ es un $A$-módulo proyectivo, por lo que $\sh{B}$ es un $\sh{O}_Y$-módulo localmente libre (\emph{loc. cit.}, n.\textsuperscript{o}~2, th. 1), y la recíproca es inmediata. Por otra parte, $f^{-1}(y)$ no es sino el espectro de la $\kappa(y)$-álgebra $\sh{B}(y)=\sh{B}_y\otimes_{\sh{O}_{Y,y}}\kappa(y)$, lo que termina de demostrar la equivalencia de $a')$ y $b)$. Para ver que $b)$ equivale a $c)$, observemos que la segunda afirmación de $b)$ equivale a que el homomorfismo $\operatorname{astr}_{\sh{B}(y)|\kappa(y)}:\sh{B}(y)\to\check{\sh{B}(y)}$ sea biyectivo; como $\sh{B}(y)=\sh{B}_y/\mathfrak m_y\sh{B}_y$ y $\check{\sh{B}(y)}=\check{\sh{B}}_y/\mathfrak m_y\check{\sh{B}}_y$ y $\sh{B}_y$ y $\check{\sh{B}}_y$ son $\sh{O}_{Y,y}$-módulos libres, se sigue de (18.2.2.4) y de Bourbaki, \emph{Alg. comm.}, cap. II, \textsection~3, n.\textsuperscript{o}~3, cor. de la prop. 6, que el homomorfismo $\operatorname{astr}_{\sh{B}_y|\sh{O}_{Y,y}}:\sh{B}_y\to\check{\sh{B}}_y$ es también biyectivo; siendo evidente la recíproca, esto termina la demostración.

Cuando una $\sh{O}_X$-álgebra $\sh{B}$ satisface las condiciones equivalentes $b)$ y $c)$ de (18.2.3), se dice que $\sh{B}$ es una $\sh{O}_X$-álgebra \emph{finita y étale}. Cuando $X=\Spec(A)$ es afín y se tiene, por tanto, $\sh{B}=\widetilde B$, donde $B$ es una $A$-álgebra, en virtud de (18.2.3) equivale decir que $\sh{B}$ es una $\sh{O}_X$-álgebra finita y étale o que $B$ es una $A$-álgebra finita y étale (en el sentido de (17.3.2)).

\begin{corollary}[18.2.4]
\label{IV.18.2.4}
Sea $f:X\to Y$ un morfismo finito y de presentación finita, y pongamos $\sh{B}=f_*(\sh{O}_X)$. Sea $y$ un punto de $Y$. Las condiciones siguientes son equivalentes:
\begin{enumerate}
  \item[a)] Existe un entorno abierto $U$ de $y$ en $Y$ tal que la restricción $f^{-1}(U)\to U$ de $f$ sea un morfismo étale.
  \item[b)] $\sh{B}_y$ es un $\sh{O}_{Y,y}$-módulo libre de tipo finito y $\sh{B}_y\otimes_{\sh{O}_{Y,y}}\kappa(y)$ es una $\kappa(y)$-álgebra separable.
\end{enumerate}
\end{corollary}

Es claro que $a)$ implica $b)$ en virtud de (18.2.3). Por otra parte, $\sh{B}$ es un $\sh{O}_Y$-módulo de presentación finita (1.6.3 y 1.4.7), de modo que, si $\sh{B}_y$ es un $\sh{O}_{Y,y}$-módulo libre, existe un entorno abierto $U$ de $y$ en $Y$ tal que $\sh{B}|U$ sea un $\sh{O}_U$-módulo localmente libre ($0_{\mathrm I}$, 5.2.7); además, por hipótesis, como el homomorfismo $\operatorname{astr}_{\sh{B}_y|\sh{O}_{Y,y}}:\sh{B}_y\to\check{\sh{B}}_y$ es biyectivo, se sigue también de ($0_{\mathrm I}$, 5.2.7) que se puede suponer elegido $U$ de modo que

\oldpage[IV-4]{113}

el homomorfismo $\operatorname{astr}_{\sh{B}|U/\sh{O}_U}$ sea biyectivo. El hecho de que $b)$ implique $a)$ se sigue entonces de (18.2.3).

\begin{corollary}[18.2.5]
\label{IV.18.2.5}
Sean $Y$ un preesquema cuasicompacto o localmente noetheriano y $f:X\to Y$ un morfismo finito y de presentación finita; pongamos $\sh{B}=f_*(\sh{O}_X)$. Supongamos que, para todo punto cerrado $y$ de $Y$, $\sh{B}_y$ sea un $\sh{O}_{Y,y}$-módulo libre y $\sh{B}_y\otimes_{\sh{O}_{Y,y}}\kappa(y)$ una $\kappa(y)$-álgebra separable. Entonces $f$ es étale.
\end{corollary}

En efecto, se sigue de (18.2.4) que todo punto cerrado $y$ de $Y$ posee un entorno abierto $U$ tal que la restricción $f^{-1}(U)\to U$ de $f$ sea étale; la conclusión se sigue de que, en los dos casos considerados, todo subconjunto cerrado no vacío de $Y$ contiene un punto cerrado (5.1.11 y $0_{\mathrm I}$, 2.1.3).

\begin{corollary}[18.2.6]
\label{IV.18.2.6}
Si $f:X\to Y$ es un morfismo finito y étale, y se pone $\sh{B}=f_*(\sh{O}_X)$, entonces el $\sh{O}_Y$-homomorfismo $\operatorname{Tr}_{\sh{B}|\sh{O}_Y}:\sh{B}\to\sh{O}_Y$ (18.2.2) (denotado también por $\operatorname{Tr}_f$) es sobreyectivo.
\end{corollary}

Como la cuestión es local, en virtud de (18.2.3) se puede suponer que $Y=\Spec(A)$, $X=\Spec(B)$ con $\sh{B}=\widetilde B$, siendo $B$ un $A$-módulo libre; como en virtud de (18.2.3) la forma bilineal (18.2.1.1) es no degenerada, esto implica, en particular, que la forma lineal $\operatorname{Tr}_{B/A}$ es sobreyectiva.

\begin{remarks}[18.2.7]
\label{IV.18.2.7}
\begin{enumerate}
  \item[(i)] Cuando $f:X\to Y$ es un morfismo finito tal que $f_*(\sh{O}_X)$ sea un $\sh{O}_Y$-módulo localmente libre (resp. localmente libre de rango $n$), se dice también que $f$ es un morfismo \emph{finito localmente libre} (resp. \emph{localmente libre de rango $n$}). Esta condición, en virtud de (18.2.3), se satisface si $f$ es un morfismo finito étale, pero no implica por sí sola que $f$ sea étale, como muestra el ejemplo en que $X=\Spec(K)$ e $Y=\Spec(k)$ son espectros de cuerpos, siendo $K$ una extensión finita no separable de $k$. Cuando $f:X\to Y$ es un morfismo finito étale, se dice también que $X$ es un \emph{recubrimiento étale} de $Y$.

Obsérvese que, en este caso, $f$ es universalmente abierto y universalmente cerrado y, en particular, $f(X)$ es un subconjunto a la vez abierto y cerrado de $Y$.

Se dice que un recubrimiento étale $X$ de $Y$ es \emph{trivial} si $X$ es \emph{suma} de un número finito de preesquemas \emph{isomorfos} a $Y$. Se dice que un recubrimiento étale $X$ de $Y$ es \emph{localmente trivial} si el morfismo $f:X\to Y$ es tal que todo punto $y\in Y$ posee un entorno abierto $U$ para el cual el recubrimiento $f^{-1}(U)$ de $U$ es trivial.

  \item[(ii)] Sea $f:X\to Y$ un morfismo finito, localmente libre de rango $n$; pongamos $\sh{B}=f_*(\sh{O}_X)$ y sea $u=\operatorname{astr}_{\sh{B}|\sh{O}_Y}:\sh{B}\to\check{\sh{B}}$; de él se deduce un homomorfismo sobre la $n$-ésima potencia exterior
\[
  \bigwedge^n u:\bigwedge^n\sh{B}\longrightarrow\bigwedge^n\check{\sh{B}}=(\bigwedge^n\sh{B})^\vee
\]
de $\sh{O}_Y$-módulos \emph{invertibles} y, por consiguiente, ($0_{\mathrm I}$, 5.4.2), un elemento
\begin{equation}
  d_{X/Y}\in\Gamma\left(Y,(\bigwedge^n\check{\sh{B}})\otimes_{\sh{O}_Y}(\bigwedge^n\check{\sh{B}})\right)
  \tag{18.2.7.1}\label{IV.18.2.7.1}
\end{equation}
que se llama el \emph{discriminante de $X$ sobre $Y$}. Además, como $(\bigwedge^n\check{\sh{B}})\otimes_{\sh{O}_Y}(\bigwedge^n\check{\sh{B}})$ es el dual de $(\bigwedge^n\sh{B})\otimes_{\sh{O}_Y}(\bigwedge^n\sh{B})$, $d_{X/Y}$ puede identificarse también con un homomorfismo
\begin{equation}
  (\bigwedge^n\sh{B})\otimes_{\sh{O}_Y}(\bigwedge^n\sh{B})\longrightarrow\sh{O}_Y
  \tag{18.2.7.2}\label{IV.18.2.7.2}
\end{equation}
y se denota por $\mathscr{D}_{X/Y}$ el ideal cuasicoherente de tipo finito de $\sh{O}_Y$, \emph{imagen} del homomorfismo (18.2.7.2), llamado también \emph{ideal discriminante de $X$ sobre $Y$}.
\end{enumerate}
\end{remarks}

\oldpage[IV-4]{114}

Dicho esto, para que el homomorfismo $u$ sea biyectivo, es necesario y suficiente que $\bigwedge^n u$ sea biyectivo, o también que la sección $d_{X/Y}$ tenga un germen \emph{invertible} en todo punto $y\in Y$, lo que se escribe asimismo $d_{X/Y}(y)\neq 0$ para todo $y$ ($0_{\mathrm I}$, 5.5.2). Equivale también decir que el ideal discriminante $\mathscr{D}_{X/Y}$ es \emph{igual a $\sh{O}_Y$}.

La terminología de «recubrimiento» introducida en (18.2.7, (i)) se justifica por la proposición siguiente:

\begin{proposition}[18.2.8]
\label{IV.18.2.8}
Sea $f:X\to Y$ un morfismo étale, separado y de tipo finito y, para todo $y\in Y$, sea $n(y)$ el número geométrico de puntos de $f^{-1}(y)$. Entonces la función $y\mapsto n(y)$ es semicontinua inferiormente en $Y$. Para que sea continua en un punto $y$ (y, por tanto, constante en un entorno de $y$), es necesario y suficiente que exista un entorno abierto $U$ de $y$ tal que la restricción $f^{-1}(U)\to U$ de $f$ sea un morfismo finito (étale).
\end{proposition}

Como $f$ es cuasifinito (17.6.1) y localmente de presentación finita, decir que $f$ es finito equivale a decir que $f$ es propio (8.11.1); además, cada fibra $f^{-1}(y)$ es geométricamente reducida sobre $\kappa(y)$. Las conclusiones se siguen, por tanto, de (15.5.9, (i) y (ii)) y de que $f$ es plano.

\begin{corollary}[18.2.9]
\label{IV.18.2.9}
Sean $Y$ un preesquema conexo y $f:X\to Y$ un morfismo étale, separado y de tipo finito. Para que $f$ sea finito (en otras palabras, para que $X$ sea un recubrimiento étale de $Y$), es necesario y suficiente que todas las fibras de $f$ tengan el mismo número geométrico de puntos.
\end{corollary}

\begin{remarks}[18.2.10]
\label{IV.18.2.10}
\begin{enumerate}
  \item[(i)] El ejemplo de la «recta afín con un punto desdoblado» (I, 5.5.11) muestra que un morfismo étale de tipo finito de preesquemas noetherianos puede ser \emph{no separado}; para este ejemplo, la primera afirmación de (18.2.8) ya no es exacta.

  \item[(ii)] Para que un morfismo separado, de tipo finito y étale $f:X\to Y$ haga de $X$ un recubrimiento \emph{localmente trivial}, es necesario y suficiente que, para todo $x\in X$, exista un entorno abierto $U$ de $y=f(x)$ y una $U$-sección $g$ de $f^{-1}(U)$ tal que $g(y)=x$. En efecto, la condición es evidentemente necesaria; que sea suficiente se sigue de que toda fibra $f^{-1}(y)$ es finita (17.6.1) y (I, 6.2.2), de la caracterización de las secciones de un $Y$-preesquema étale (17.9.3) y de la prop. (18.2.8).
\end{enumerate}
\end{remarks}
